\documentclass[11pt]{article}
\usepackage[a4paper, margin=2.0cm]{geometry}
\usepackage[T1]{fontenc}
\usepackage[utf8]{inputenc}
\usepackage{booktabs, array, amsmath, xcolor, enumitem}
\usepackage[colorlinks=true, linkcolor=blue!50!black, urlcolor=blue!50!black]{hyperref}
\setlength{\parskip}{4pt}
\sloppy

% Every number below comes from experiments/helib/pilot_tables/ (generated by
% experiments/helib/make_helib_pilot_tables.py from aggregate.py
% --scenario=helib_pilot and the recorded HElib parameters).
\newcommand{\tab}[1]{../../experiments/helib/pilot_tables/#1}
\input{\tab{numbers.tex}}

\title{HElib latency pilot\\[2pt]
\large Parameters, operation boundaries, CKKS precision and latency}
\author{Boungo Tr\'esor}
\date{\today}

\begin{document}
\maketitle

\noindent\textbf{Purpose.} As agreed (7 and 18 November): one latency pilot for HElib, documenting the security parameters actually used, the operation boundaries and the CKKS precision achieved, before any expansion of scope. HElib~v2.3.0 (commit \texttt{3e337a6}; its version header still reports 2.2.0), configuration $N{=}8192$, category~1, CKKS and BGV (HElib has no BFV), measured in one session together with SEAL, with the protocol of the OpenFHE pilot. Code: commit \texttt{\hlCommit}; raw data: \texttt{experiments/results/raw/helib\_pilot/}.

\paragraph{Summary.}
\begin{itemize}[leftmargin=*, itemsep=2pt]
  \item \textbf{The session is complete and valid}: \hlCells\ cells (\hlHelibCells\ HElib, \hlSealCells\ SEAL), \hlReps\ measured repetitions each after \hlWarm\ warm-up, all \hlAllChecks\ correctness checks passed (\hlHelibChecks\ of them HElib's); on mains, governor \hlGovernor, boost \hlBoost, \hlDate, \hlStart--\hlEnd.
  \item \textbf{HElib does not take a list of primes}: it builds its own chain from a requested bit size. The configuration was chosen by a rule stated before the session, from \hlProbeCandidates\ probed candidates. Both HElib configurations are within the HE standard's bound for $N{=}8192$ at 128~bits, but they are not SEAL's: CKKS $Q = \hlCkksQ$, $P = \hlCkksP$ bits; BGV $Q = \hlBgvQ$, $P = \hlBgvP$ bits (SEAL: $Q = \hlSealQ$, $P = \hlSealPbits$ bits).
  \item \textbf{No HElib operation has SEAL's boundary}, so no ratio between the libraries is formed; the values are reported side by side.
  \item \textbf{CKKS precision}: \hlPrecHelibMin--\hlPrecHelibMax\ bits for HElib, \hlPrecSealMin--\hlPrecSealMax\ bits for SEAL (same session, same error measure). With HElib's two ciphertext primes, a higher precision setting leaves depth~\hlProbeHighPrecDepth.
  \item \textbf{Session conditions}: SEAL's latency in this session agrees with the archived SEAL campaign within $\pm$10\% (TOST) for \hlReproEq\ of its \hlReproN\ cells.
\end{itemize}

\section{Parameters actually used}

HElib is configured by the cyclotomic index $m = 2N$, the bit size of the ciphertext modulus chain (\texttt{bits}), the number of key-switching columns ($c$) and the CKKS precision or the BGV plaintext modulus $p$. It then chooses three sets of primes: ciphertext primes (the modulus $Q$ of a fresh ciphertext), special primes ($P$, used only in key switching) and ``small'' primes (\hlCkksSmallN\ here), which only lower a ciphertext's modulus in finer steps. HElib's source shows that its keys are built over the ciphertext and special primes only (\texttt{fullPrimes()}), so security depends on $\log_2 QP$ over these two sets, which is also what HElib's own estimate uses.

A probe program (\texttt{experiments/helib/probe/}) recorded, for each candidate, the primes HElib builds, $\log_2 QP$, HElib's security estimate, the depth reached (fresh inputs, multiplication and relinearization repeated, every result decrypted and checked) and the CKKS precision. The rule (\texttt{make\_params.py}): $c$ equal to the number of ciphertext primes (one key-switching digit per prime, as SEAL), at least the grid's depth (2), $\log_2 QP \le 218$, HElib's estimate $\ge 128$, and the smallest $\log_2 QP$. The harness compares the primes HElib builds with the recorded ones and stops if they differ.

\begin{center}\footnotesize\setlength{\tabcolsep}{4pt}\input{\tab{params.tex}}\end{center}

Two findings of the probe:
\begin{itemize}[leftmargin=*, itemsep=2pt]
  \item \textbf{HElib's \texttt{SecKey} overrides \texttt{Encrypt} with secret-key encryption}, also when it is called through a public-key reference. Secret-key encryption adds less noise; a first probe that used it reported BGV depth~2 with two ciphertext primes. With public-key encryption on a separate \texttt{PubKey} object (used everywhere in the pilot), two primes give BGV depth \hlBgvTwoPrimesDepth, and the next candidate has three primes and depth \hlBgvDepth: no BGV configuration has exactly the grid's depth.
  \item \textbf{CKKS precision and depth compete}: with $Q = \hlCkksQ$, HElib's default precision setting (\hlCkksPrecParam) gives depth~\hlCkksDepth; settings up to \hlProbeHighPrecMax\ give more precision but depth~\hlProbeHighPrecDepth. The pilot keeps the default, for the grid's depth.
\end{itemize}

\section{Operation boundaries}

Each HElib operation is timed through HElib's own API, with inputs prepared before the timer and every result checked after it (\texttt{experiments/helib/src/bench\_helib.cpp}).
\begin{itemize}[leftmargin=*, itemsep=1pt]
  \item \textbf{Key generation}: context construction (prime chain, algebra, tables), secret key, public key and the relinearization key $s^2 \to s$ only (\texttt{GenSecKey(0, 2)}; HElib's default also builds $s^3 \to s$, which SEAL does not).
  \item \textbf{Encryption}: public-key encryption of a plaintext encoded before the timer, including HElib's conversion of the encoding to its internal representation.
  \item \textbf{Decryption}: to the plaintext polynomial, before decoding; it includes HElib's check of the noise bound and the reconstruction of integer coefficients (SEAL's CKKS decryption stops before that). \textbf{Decryption with decoding}: the application-level call, an end-to-end API cost.
  \item \textbf{Addition}$^\ast$: in place (\texttt{+=}), without the copy of a new result (SEAL's \texttt{add} writes a new ciphertext).
  \item \textbf{Multiplication}: \texttt{multLowLvl}, without relinearization; HElib first switches the inputs to the modulus it computes for the product, then forms the tensor product. SEAL's multiplication is the tensor product alone.
  \item \textbf{Relinearization}: \texttt{reLinearize} on a prepared product (drops small and special primes, then key switching).
  \item \textbf{Rotation}: by one position along the first dimension, with one key-switching key (for that generator only, as SEAL's single Galois key); for BGV, each of the two rows of $N/2$ slots rotates cyclically, as in SEAL.
\end{itemize}

\section{Results}

Protocol: \hlReps\ measured repetitions after \hlWarm\ warm-up, inner loop \hlInner\ calls per repetition (one call for key generation, rotation and BGV decryption with decoding), randomized order, one process per cell pinned to one core; SEAL cells with the same settings. The inner loop is smaller than in the SEAL and OpenFHE campaigns: every call needs its own fresh HElib input, which takes about 30~ms for BGV outside the timing. Mean latency (ms):

\begin{center}\footnotesize\setlength{\tabcolsep}{6pt}\input{\tab{latency.tex}}\end{center}
{\footnotesize $^\ast$ HElib: in place; SEAL: into a new ciphertext. CKKS precision: $-\log_2$ of the largest normalized error over the repetitions of each operation, range over the operations.}

\begin{itemize}[leftmargin=*, itemsep=2pt]
  \item \textbf{BGV key generation and decoding are slow in HElib} (\hlLatHElibBgvKeygen\ and \hlLatHElibBgvDecryptFull~ms, against \hlLatHElibCkksKeygen\ and \hlLatHElibCkksDecryptFull~ms for CKKS). A plausible cause, not measured here, is the construction and use of HElib's slot structure for $p = \hlBgvPlainMod$; the parts of these operations were not timed separately.
  \item \textbf{HElib's BGV multiplication} (\hlLatHElibBgvMultiply~ms, against \hlLatHElibCkksMultiply~ms for CKKS) includes HElib's automatic modulus switching, which SEAL's multiplication does not perform.
  \item \textbf{HElib's own warnings} are kept with the data (\texttt{helib.log}): \hlLogKs\ ``KS-noise-ratio'' warnings, all from the two rotation cells (HElib reports that the key switch of a fresh ciphertext at full modulus adds much more noise than the ciphertext held; every rotated result still decrypted correctly), and \hlLogTarget\ notices about prime sizes when a BGV context is built; no other warning (\hlLogOther).
  \item \textbf{SEAL in this session} agrees with its archived campaign within $\pm$10\% for \hlReproEq\ of \hlReproN\ cells; the others are \hlReproOut. The smaller inner loop may contribute; this was not tested.
\end{itemize}

\section{Scope and next step}

This is a pilot: one configuration, one session, latency only. Its values describe HElib~v2.3.0 under the stated parameters and boundaries; they are not a comparison of the libraries' costs for the same computation. Extending HElib to the grid would require the same probe for every configuration and, for a ratio with SEAL, operations with matching boundaries. The Lattigo pilot follows the same procedure.

\end{document}
